% Copyright 2012-2014 by abnTeX2 group at http://abntex2.googlecode.com/ 
%% %% Copyright 2023 Matheus Miceli e Sávio Luis - Centro de Recursos Computacionais - CRC/ICEI PUC Minas
% --------------------------------------------------------------------------------------------------------
% abnTeX2: Modelo de Artigo Acadêmico em conformidade com ABNT NBR 6022:2003 Modelo Poster - Versão 1.0
% --------------------------------------------------------------------------------------------------------

\documentclass[article,9pt,a4paper]{abntex2}

\setlength{\headheight}{5pt} % decrease the value of \headheight
\setlength{\headsep}{5pt} % decrease the value of \headsep
\setlength{\footskip}{5pt} % decrease the value of \footskip


% ---------------------------
% Pacotes fundamentais ------
% ---------------------------
\usepackage{caption}
\usepackage{multicol} % Colunas
\setlength{\columnsep}{1cm} % espaçamento colunas
\usepackage{lmodern}			% Usa a fonte Latin Modern
\usepackage[T1]{fontenc}		% Selecao de codigos de fonte.
\usepackage[utf8]{inputenc}		% Codificacao do documento (conversão automática dos acentos)
\usepackage{indentfirst}		% Indenta o primeiro parágrafo de cada seção.
\usepackage{nomencl} 			% Lista de simbolos
\usepackage{color}				% Controle das cores
\usepackage{graphicx}
\graphicspath{{images/}}
\usepackage{microtype} 			% para melhorias de justificação
% Pacotes adicionais, usados apenas no âmbito do Modelo Canônico do abnteX2
% ---
\usepackage{lipsum}				% para geração de dummy text
% ------------------------
% Pacotes de citações ----
% ------------------------
\usepackage[brazilian,hyperpageref]{backref}	 % Paginas com as citações na bibl
	
% informações do PDF
\makeatletter
\hypersetup{
     	%pagebackref=true,
		pdftitle={\@title}, 
		pdfauthor={\@author},
    	pdfsubject={Modelo de artigo científico com abnTeX2},
	    pdfcreator={LaTeX with abnTeX2},
		pdfkeywords={abnt}{latex}{abntex}{abntex2}{atigo científico}, 
		colorlinks=true,       		% false: boxed links; true: colored links
    	linkcolor=blue,          	% color of internal links
    	citecolor=blue,        		% color of links to bibliography
    	filecolor=magenta,      		% color of file links
		urlcolor=blue,
		bookmarksdepth=4
}
\makeatother
% --- 

% -------------------------------------------------
% CONFIGURAÇÃO DAS MARGENS ------------------------
% -------------------------------------------------
\setlrmarginsandblock{0.7cm}{0.7cm}{*} % ESPAÇO NA HORIZONTAL
\setulmarginsandblock{0.5cm}{0.5cm}{*} % ESPAÇO NA VERTICAL
\checkandfixthelayout

% ------------------------------------------------------
% CONFIGURAÇÃO DOS ESPAÇOS ENTRE LINHAS E PARÁGRAFOS ---
% ------------------------------------------------------
% TAMANHO DO PARÁGRAFO
\setlength{\parindent}{0.5cm}
% ESPAÇO ENTRE UM PARÁGRAFO E OUTRO
\setlength{\parskip}{\onelineskip}  % tente também \onelineskip
% ESPAÇAMENTO SIMPLES
\SingleSpacing
% NUMERAÇÃO DAS PÁGINAS ------
\pagenumbering{gobble}

% -------------------------------------------------
% INÍCIO DO DOCUMENTO -----------------------------
% -------------------------------------------------
\begin{document}


% ---------------
% CABEÇALHO -----
% ---------------

\begin{figure}[!h]
    \begin{minipage}{0.13\linewidth}
    \includegraphics[scale=.17]{images/logoPUC.pdf}    
    \end{minipage}
    \begin{minipage}{0.8\linewidth}
         \textbf{Pontifícia Universidade Católica de Minas Gerais} \\
         Instituto de Ciências Exatas de Informática (ICEI)
    \end{minipage}
\end{figure}

\vspace{-1cm}

% --- TÍTULO 
\begin{center}
    \huge \textbf{Título Do Trabalho} \\ [0.3cm]
    \large\textit{Trabalho Conclusão De Cursos I - Junho 2018}\\
    \large \textit{Nome Do Aluno \footnote{Nome do Aluno e Breve Descrição} \& Nome Do Orientador\footnote{Nome do Orientador e Breve Descrição}}\\
    \large {aluno@sga.pucminas.br, professor@pucminas.br}
\end{center}

\vspace{0.3cm}
% FIM CABEÇALHO

% ----------------------------------------------------------
% Introdução
% ----------------------------------------------------------
\begin{multicols}{2}

\section {Introdução}
Apresentar o contexto do trabalho. Fazer isso com duas ou quatro frases, demarcando seu trabalho e sua importância na Ciência da Computação.

\section{Descrição do problema}

\subsection{O Problema}
Abordar o problema escolhido para o trabalho de forma clara e objetiva. O problema poderá ser apresentado em forma de pergunta.

\subsection{Objetivos}
Apresentar o objetivo geral do trabalho e os objetivos específicos, caso existam.

\subsection{Justificativa}
Apresentar a justificativa do desenvolvimento de seu trabalho tendo em vista a importância do tema em Ciência da Computação e a contribuição do mesmo.(letra Calibri 40 - O fundo poderá ser branco - preferencialmente - ou colorido, à sua escolha).

\begin{center}    
\includegraphics[width=5cm]{placeholder.png}
\captionof{figure}{Titulo e origem}
\end{center}


\section{Referencial Teórico}
(Apresente as ideias dos principais autores que você leu. Estas ideais deverão ser a base para sustentar a metodologia escolhida e as conclusões. Este item poderá ser subdivido à vontade. (letra Calibri 40).Lorem ipsum dolor sit amet, consectetur adipiscing elit. Phasellus et mi feugiat, rutrum mi vitae, rutrum ante. Sed quis ornare sem. Etiam mattis ut dui ac consequat. Nulla iaculis urna et sagittis viverra. Aenean fringilla nisi et nisi tincidunt molestie.Lorem ipsum dolor sit amet, consectetur adipiscing elit. Phasellus et mi feugiat, rutrum mi vitae, rutrum ante. Sed quis ornare sem. Etiam mattis ut dui ac consequat. Nulla iaculis urna et sagittis viverra. Aenean fringilla nisi et nisi tincidunt molestie.Lorem ipsum dolor sit amet, consectetur adipiscing elit. Phasellus et mi feugiat, rutrum mi vitae, rutrum ante. Sed quis ornare sem. Etiam mattis ut dui ac consequat. Nulla iaculis urna et sagittis viverra. Aenean fringilla nisi et nisi tincidunt molestie.


\begin{equation}
E = mc^{2}
\label{eqn:Einstein}
\end{equation}

Lorem ipsum dolor sit amet, consectetur adipiscing elit. Phasellus et mi feugiat, rutrum mi vitae, rutrum ante. Sed quis ornare sem. Etiam mattis ut dui ac consequat. Nulla iaculis urna et sagittis viverra. Aenean fringilla nisi et nisi tincidunt molestie. Pellentesque nunc ipsum, elementum eget pretium id, lobortis at nisi. Etiam non arcu vulputate, mollis diam vel, consectetur lectus. Quisque semper mollis est a porta.

Praesent et leo condimentum magna facilisis sodales ut in mi. Suspendisse a ultrices ligula. Nullam molestie eros non volutpat porttitor. Suspendisse ac hendrerit tellus, nec ultricies ligula. Vivamus congue nulla et pretium consequat. Vestibulum tempus neque nec tincidunt molestie. Aliquam ornare, sem ac egestas venenatis, purus diam scelerisque felis, vel convallis sem velit sed arcu.

\begin{center}
\begin{tabular}{l l l}
\toprule
\textbf{Treatments} & \textbf{Response 1} & \textbf{Response 2}\\
\midrule
Treatment 1 & 0.0003262 & 0.562 \\
Treatment 2 & 0.0015681 & 0.910 \\
Treatment 3 & 0.0009271 & 0.296 \\
\bottomrule
\end{tabular}
\captionof{table}{Table caption}
\end{center}

\section{Metodologia}
Detalhamento de como foi ou de como será realizado o trabalho. Procure qualificar a metodologia (Pesquisa Exploratória, Pesquisa Experimental; Estudo de Caso, etc.. ). Em seguida liste os passos para a consecução do trabalho. Utilize o Smart Art - no Power Point – Inserir – Smart Art) para correlacionar ideias e apresentar esquemas cognitivos, caso assim deseje. Vide exemplo Figura 1.


\begin{eqnarray}
\cos\bar{\phi}_k Q_{j,k+1,t} + Q_{j,k+1,x}+\frac{\sin^2\bar{\phi}_k}{T\cos\bar{\phi}_k} Q_{j,k+1} &=&\nonumber\\ 
-\cos\phi_k Q_{j,k,t} + Q_{j,k,x}-\frac{\sin^2\phi_k}{T\cos\phi_k} Q_{j,k}\label{edgek}
\end{eqnarray}


\section{Próximos passos (TCC II)}
O manual do TEX \cite{audy2009} pode ser usado para aprendê-lo, e o livro do Lamport \cite{farinelli2007} para aprender o LATEX, mas se quiser ir a fundo tem que ver como o TEXquero o os parágrafos em linhas \cite{Leite2009}.

\begin{itemize}
    \item Pellentesque eget orci eros. Fusce ultricies, tellus et pellentesque fringilla,
    \item Pellentesque eget orci eros. Fusce ultricies, tellus et pellentesque fringilla,
\end{itemize}

\renewcommand\bibsection
{
   \section*{\vspace{-\baselineskip}}
   
}
\section*{Referências}\vspace{-0.8cm}
\bibliographystyle{abbrv} % Plain referencing style
\bibliography{referencias} % Use the example bibliography file %sample.bib

\section*{Agradecimentos}
\textit{\small This research was undertaken text is derived from sections 1.10.33 of Cicero’s De finibus bonorum et malorumtext is derived from
sections 1.10.33 of Cicero’s De finibus bonorum et malorumtext is derived from sections 1.10.33 of Cicero’s De finibus bonorum et
malorum.}

\end{multicols}
\end{document}
